% Attività dei rifiuti radioattivi nel tempo: da 640 kBq a 80 kBq in 30 h
% (3 emivite, T1/2 = 10 h), a 10 kBq in 60 h (6 emivite).
% Esercizio ud7-radio-rifiuti (soluzione)
\begin{tikzpicture}
\begin{axis}[
  width=0.68\linewidth, height=5.4cm,
  axis lines=left, xmin=0, xmax=66, ymin=0, ymax=720,
  xlabel={tempo $t$ (\si{\hour})}, ylabel={$\mathcal{A}$ (\si{\kilo\becquerel})},
  xtick={0,10,20,30,40,50,60}, ytick={0,80,160,320,640},
  tick label style={font=\small}, label style={font=\small},
]
  \addplot[attenzione, very thick, domain=0:64, samples=100] {640*2^(-x/10)};
  \foreach \t/\a in {10/320, 20/160, 30/80} {
    \edef\tmp{\noexpand\draw[dashed, gray] (axis cs:\t,0) -- (axis cs:\t,\a) -- (axis cs:0,\a);}\tmp
  }
  \addplot[only marks, mark=*, formula] coordinates {(0,640) (30,80) (60,10)};
  \node[formula, right, font=\small] at (axis cs:1.5,640) {prima misura};
  \node[formula, above right, font=\small, align=left] at (axis cs:30,80) {seconda misura\\dopo 3 emivite};
  \node[formula, above, font=\small, align=center] at (axis cs:58,10) {\SI{10}{\kilo\becquerel}\\si smaltisce};
\end{axis}
\end{tikzpicture}
